% ---------------------------------------------------------------------------
% All frames for the GAM-Softmax talk. Included by slides.tex and notes.tex.
% ---------------------------------------------------------------------------

% =========================================================== TITLE
\begin{frame}[plain]
  \vfill
  \centering
  {\color{gamgrey}\small GENERALIZED ADAPTIVE MARGIN SOFTMAX}\\[1.0em]
  {\color{gamblue}\huge\bfseries GAM-Softmax}\\[0.6em]
  {\color{gamblue}\large Can we make the softmax margin \emph{smart}\\ instead of a single number?}\\[1.4em]
  {\color{gamblue}\rule{0.4\textwidth}{1pt}}\\[1.2em]
  {\small Ashish Gupta}\\[0.3em]
  {\color{gamgrey}\small One RTX 3050 laptop GPU \quad$\cdot$\quad 23 May -- 31 July 2026}
  \vfill
  \note{%
    Greeting. This talk is about the last layer of a neural-network classifier
    --- the softmax --- and about one number inside a published improvement to
    it. I will tell you what softmax is and why it needs improving, what the
    published method does, what I proposed instead, what I built, what I ran,
    and what came out. I will be straight with you about the outcome: the main
    hypothesis did \emph{not} hold, and I will show you exactly how I know
    that. One sub-part did work, and that part is the interesting maths.%
  }
\end{frame}

% =========================================================== ROADMAP
\begin{frame}{Where this talk goes}
  \begin{columns}[T]
    \begin{column}{0.52\textwidth}
      \small
      \begin{enumerate}\setlength\itemsep{0.55em}
        \item \textbf{Introduction} --- what softmax is, and the flaw in how we train it
        \item \textbf{Background} --- AS-Softmax and its single number $\dlt$
        \item \textbf{Methodology} --- turning $\dlt$ into a function
        \item \textbf{Experimentation} --- what I ran and what I measured
        \item \textbf{Results} --- the numbers, and the verdict
      \end{enumerate}
    \end{column}
    \begin{column}{0.44\textwidth}
      \small
      \begin{block}{The whole talk in one line}
        A classifier is normally told \emph{``be 100\,\% sure, forever''}.
        A better instruction is \emph{``win by a margin $\dlt$, then stop''}.
        I asked whether that margin should be \alert{one number} or \alert{a function}.
      \end{block}
    \end{column}
  \end{columns}
  \note{%
    Five sections. Roughly two minutes each. The single sentence on the right is
    the thing to hold on to: cross-entropy, the standard training objective,
    tells the network to be completely certain and never lets it stop. A margin
    objective tells it to win by a fixed gap and then stop. My question was
    whether that gap should stay a single constant, or become a function that
    varies. Everything else is detail.%
  }
\end{frame}

% ===========================================================================
\section{Introduction}
% ===========================================================================

\begin{frame}{Where the softmax sits}
  \framesubtitle{The last layer of essentially every classifier}
  \centering
  \resizebox{0.97\textwidth}{!}{%
  \begin{tikzpicture}[
      box/.style={draw=gamblue, fill=gamlight, rounded corners=2pt,
                  minimum height=1.5em, inner sep=4pt, font=\scriptsize, align=center},
      arr/.style={-{Stealth[length=4pt]}, gamblue, thick}]
    \node[box] (x) {input $x$\\{\tiny a document}};
    \node[box, right=0.7cm of x] (f) {encoder $f_\theta$\\{\tiny BERT}};
    \node[box, right=0.7cm of f] (z) {logits $z\in\R^{n}$\\{\tiny $n$ real scores}};
    \node[box, right=0.7cm of z, fill=gamblue!18] (s) {\textbf{softmax}};
    \node[box, right=0.7cm of s] (p) {$p\in\R^{n}$\\{\tiny $p_i>0,\ \sum_i p_i=1$}};
    \node[box, right=0.7cm of p] (y) {prediction\\{\tiny $\arg\max_i p_i$}};
    \draw[arr] (x)--(f); \draw[arr] (f)--(z); \draw[arr] (z)--(s);
    \draw[arr] (s)--(p); \draw[arr] (p)--(y);
  \end{tikzpicture}}

  \vspace{0.6em}
  \begin{columns}[T]
    \begin{column}{0.48\textwidth}
      \small\textbf{The job it performs}\\[0.2em]
      The network emits $n$ unbounded real numbers (``logits'') --- they can be
      negative, they do not sum to anything. Softmax is the map
      \[ \R^{n}\ \longrightarrow\ \Delta^{n-1} \]
      that turns them into a \emph{probability distribution} over the $n$ classes.
    \end{column}
    \begin{column}{0.48\textwidth}
      \small\textbf{Where it is used}\\[0.2em]
      \begin{itemize}\setlength\itemsep{0.15em}
        \item every classifier: text, image, audio
        \item every language model: the distribution over the next token ($n\approx50{,}000$)
        \item attention weights inside a Transformer
        \item action probabilities in reinforcement learning
      \end{itemize}
    \end{column}
  \end{columns}
  \note{%
    Read the pipeline left to right. A document goes into an encoder --- in my
    experiments, BERT --- which produces twenty raw scores, one per class. Those
    scores are arbitrary real numbers. To train with them, and to report a
    confidence, we need a probability distribution: numbers that are positive
    and add to one. That conversion is the softmax, and it maps
    n-dimensional real space onto the probability simplex. It is not a niche
    component --- it is the output layer of every classifier, the next-token
    distribution of every language model, and the attention weights inside
    Transformers.%
  }
\end{frame}

\begin{frame}{The softmax function}
  \[
    p_i \;=\; \operatorname{softmax}(z)_i \;=\; \frac{e^{z_i}}{\displaystyle\sum_{j=1}^{n} e^{z_j}},
    \qquad i = 1,\dots,n
  \]
  \vspace{-0.4em}
  \begin{columns}[T]
    \begin{column}{0.5\textwidth}
      \small\textbf{Why the formula is built this way}
      \begin{enumerate}\setlength\itemsep{0.25em}
        \item \textbf{Positive:} $e^{z_i}>0$ always $\Rightarrow p_i>0$.
        \item \textbf{Normalised:} dividing by the sum forces $\sum_i p_i = 1$.
        \item \textbf{Order-preserving:} $e^{x}$ is strictly increasing, so
              $z_i>z_j \Leftrightarrow p_i>p_j$. The prediction is unchanged.
        \item \textbf{Smooth:} differentiable everywhere $\Rightarrow$ trainable by gradient descent.
      \end{enumerate}
    \end{column}
    \begin{column}{0.48\textwidth}
      \small\textbf{Two properties I will use later}
      \begin{itemize}\setlength\itemsep{0.25em}
        \item \alert{Shift invariance:}
              $\operatorname{softmax}(z + c\mathbf{1}) = \operatorname{softmax}(z)$,
              since $e^{z_i+c}=e^{c}e^{z_i}$ cancels top and bottom.
              \emph{Only differences of logits carry information} --- which is
              why a margin is defined on a \emph{difference}.
        \item \alert{Temperature:} $\operatorname{softmax}(z/T)$ tends to
              $\arg\max$ as $T\!\to\!0$ and to uniform as $T\!\to\!\infty$.
              Softmax is the \emph{smooth} version of $\max$ --- hence the name.
      \end{itemize}
    \end{column}
  \end{columns}
  \note{%
    The formula: exponentiate every score, then divide by the total. Each design
    choice earns its place. Exponentiation guarantees positivity. Dividing by
    the sum guarantees the numbers add to one. Because the exponential is
    strictly increasing, the ordering of the scores is preserved --- softmax
    never changes which class wins. And it is smooth, so we can differentiate
    through it, which is what makes training possible at all. Two properties
    matter later. First, shift invariance: adding the same constant to every
    logit changes nothing, so only \emph{differences} between logits are
    meaningful --- that is exactly why the margin I will show you is defined on
    a difference of probabilities. Second, with a temperature parameter softmax
    interpolates between the hard maximum and the uniform distribution: it is a
    soft max, which is where the name comes from.%
  }
\end{frame}

\begin{frame}{A worked example}
  \framesubtitle{Four classes, one document}
  \small
  \[
    \text{exponentiate}\ \longrightarrow\ \text{sum}\ \longrightarrow\ \text{divide}
    \qquad\qquad
    p_i = \frac{e^{z_i}}{\sum_j e^{z_j}}
  \]
  \vspace{-0.8em}
  \begin{columns}[T]
    \begin{column}{0.47\textwidth}
      \begin{table}
        \centering\small
        \begin{tabular}{@{}lrrr@{}}
          \toprule
          class & $z_i$ & $e^{z_i}$ & $p_i$ \\
          \midrule
          sport    & 3.2 & 24.53 & \hit{0.775} \\
          politics & 1.3 &  3.67 & 0.116 \\
          medicine & 0.2 &  1.22 & 0.039 \\
          religion & 0.8 &  2.23 & 0.070 \\
          \midrule
          \multicolumn{2}{@{}l}{total} & 31.65 & 1.000 \\
          \bottomrule
        \end{tabular}
      \end{table}
    \end{column}
    \begin{column}{0.5\textwidth}
      \textbf{Check the shift invariance.} Subtract $3.2$ from every logit:
      \[ z' = (0,\,-1.9,\,-3.0,\,-2.4) \]
      \[ e^{z'} = (1,\ 0.150,\ 0.050,\ 0.091),\ \ \text{sum}=1.291 \]
      \[ p' = (0.775,\ 0.116,\ 0.039,\ 0.070) = p \]
      \takeaway{Identical. The absolute size of a logit means nothing; only the
      gaps between logits do. \textbf{This is the whole reason a ``margin'' is
      the natural thing to control.}}
    \end{column}
  \end{columns}
  \note{%
    Concretely, with four classes. Exponentiate: 3.2 becomes 24.5, and note how
    violently the exponential separates things --- a logit gap of 1.9 became a
    factor of about seven. Sum the exponentials, divide, and you get a valid
    distribution: 77.5 percent on sport. Now the check on the right. Subtract
    3.2 from every logit so the largest becomes zero. The exponentials all
    change, the sum changes, but every probability is exactly what it was. That
    is shift invariance, and it is the mathematical reason the rest of this talk
    is about \emph{differences}: only the gap between the correct class and a
    competitor carries information.%
  }
\end{frame}

\begin{frame}{How we train it: cross-entropy}
  \framesubtitle{And the flaw hiding in its gradient}
  \small
  With true class $t$ and one-hot label $y$, the loss is
  \[
    L \;=\; -\sum_{i=1}^{n} y_i \log p_i \;=\; \alert{-\log p_t}
    \qquad\text{(only the true-class term survives)}
  \]
  Differentiating through the softmax gives a famously clean gradient:
  \[
    \frac{\partial L}{\partial z_j} \;=\; p_j - y_j
    \qquad\Longrightarrow\qquad
    \frac{\partial L}{\partial z_t} \;=\; \mhit{p_t - 1},
    \qquad
    \frac{\partial L}{\partial z_j} \;=\; p_j \ \ (j\neq t)
  \]

  \begin{columns}[T]
    \begin{column}{0.54\textwidth}
      \textbf{Read the gradient.}
      \begin{itemize}\setlength\itemsep{0.2em}
        \item It vanishes \emph{only} when $p_t = 1$ exactly and every $p_j = 0$.
        \item But $p_i = e^{z_i}/\sum_j e^{z_j} > 0$ strictly --- $p_t=1$ is
              \alert{unreachable} for finite logits.
      \end{itemize}
    \end{column}
    \begin{column}{0.42\textwidth}
      \takeaway{\textbf{Cross-entropy never stops pushing.} At $p_t = 0.99$ it
      still pulls with $-0.01$, and it will keep pulling for as long as you
      train. There is no notion of ``this example is done''.}
    \end{column}
  \end{columns}
  \note{%
    Training uses cross-entropy. Because the label is one-hot, the sum collapses
    to a single term: minus the log of the probability assigned to the correct
    class. Differentiate that through the softmax and you get one of the
    tidiest results in machine learning --- the gradient with respect to the
    logits is simply predicted minus actual, p minus y. Now look at when that
    gradient is zero. It requires the correct class to have probability exactly
    one. But softmax output is a ratio of positive exponentials, so every
    probability is strictly positive and exactly one is unreachable. The
    consequence is the flaw: cross-entropy has no concept of an example being
    finished. At ninety-nine percent confidence it is still pulling, and it will
    pull forever.%
  }
\end{frame}

\begin{frame}{Why ``never stops'' is a real problem}
  \begin{columns}[T]
    \begin{column}{0.55\textwidth}
      \small
      \begin{enumerate}\setlength\itemsep{0.5em}
        \item \textbf{Wasted effort.} With $n=20$ classes, 19 of the 20 slots
              per example get gradient at every step --- including classes that
              were hopeless from epoch one.
        \item \textbf{Overconfidence.} The objective's only fixed point is total
              certainty, so probabilities get pushed to extremes and stop being
              honest confidences.
        \item \alert{Memorisation of wrong labels.} If a training label is
              \emph{incorrect}, cross-entropy pushes the model to fit it just as
              hard. Given enough steps it succeeds.
      \end{enumerate}
    \end{column}
    \begin{column}{0.42\textwidth}
      \small
      \begin{alertblock}{The point that drives this project}
        A deep network has enough capacity to fit \emph{arbitrary} labels,
        including random ones. Cross-entropy gives it no reason to stop --- so
        with noisy labels, it memorises the noise.
      \end{alertblock}
      \takeaway{The natural fix: \textbf{stop pushing once the answer is clearly
      right}. Now we need to define ``clearly''.}
    \end{column}
  \end{columns}
  \note{%
    Three consequences. First, wasted effort: with twenty classes, nineteen
    slots per example receive gradient every single step, including classes
    that were never plausible. Second, overconfidence --- the only fixed point
    of the objective is total certainty, so the probabilities stop being
    trustworthy confidences. Third, and this is the one that drives the whole
    project: if a training label is simply wrong, cross-entropy pushes the model
    to fit the wrong label with exactly the same force. And a large network has
    enough capacity to succeed --- it will memorise the mistake. So the obvious
    fix is to give the objective a way to say ``this example is done, stop''.
    The rest of the talk is about how to define \emph{done}.%
  }
\end{frame}

% ===========================================================================
\section{Background}
% ===========================================================================

\begin{frame}{AS-Softmax (Lv et al., 2023) --- the method I build on}
  \small
  \textbf{The rule.} For a training example with true class $t$, \emph{drop}
  a competing class $j$ from the loss as soon as it is beaten by a margin $\dlt$:
  \vspace{-0.4em}
  \[
    \text{mask } j \iff p_t - p_j \geq \dlt,
    \qquad
    \keep(x) = \{t\} \cup \{\, j\neq t : p_t - p_j < \dlt \,\},
    \qquad
    L_{\mathrm{AS}} = -\log \frac{e^{z_t}}{\sum_{k \in \keep(x)} e^{z_k}}
  \]
  \vspace{-1.0em}

  \begin{columns}[T]
    \begin{column}{0.5\textwidth}
      \begin{center}
      \begin{tikzpicture}[scale=0.78]
        \def\bw{0.6}
        \draw[gamgrey, -{Stealth[length=4pt]}] (-0.2,0) -- (5.0,0);
        \draw[gamgrey, -{Stealth[length=4pt]}] (-0.2,0) -- (-0.2,2.3) node[above, font=\tiny] {$p$};
        \fill[gamblue]   (0.2,0) rectangle ++(\bw,2.0);
        \fill[gamaccent] (1.4,0) rectangle ++(\bw,1.6);
        \fill[gamgrey]   (2.6,0) rectangle ++(\bw,0.28);
        \fill[gamgrey]   (3.8,0) rectangle ++(\bw,0.12);
        \node[font=\tiny, below] at (0.5,0)  {$t$ \;0.50};
        \node[font=\tiny, below] at (1.7,0)  {$j_1$ 0.40};
        \node[font=\tiny, below] at (2.9,0)  {$j_2$ 0.07};
        \node[font=\tiny, below] at (4.1,0)  {$j_3$ 0.03};
        \draw[dashed, gamblue, thick] (0.1,1.4) -- (4.6,1.4);
        \node[font=\tiny, gamblue, above] at (3.4,1.42) {$p_t-\dlt = 0.35$};
        \node[font=\tiny, gamaccent] at (1.7,1.85) {kept};
        \node[font=\tiny, gamgrey] at (2.9,0.55) {masked};
        \node[font=\tiny, gamgrey] at (4.1,0.42) {masked};
      \end{tikzpicture}
      \end{center}
    \end{column}
    \begin{column}{0.47\textwidth}
      \scriptsize
      With $\dlt = 0.15$: $j_1$ is only $0.10$ behind, so it is \emph{still a
      threat} and stays in the loss. $j_2, j_3$ are $0.43$ and $0.47$ behind ---
      beaten by more than $\dlt$, so they are dropped and contribute no gradient.
      \vspace{0.2em}
      \takeaway{The instruction changes from \emph{``beat everything by
      everything''} to \emph{``beat everything by $\dlt$, then stop''}.
      Reported: better accuracy, better calibration, $\approx 1.2\times$ faster.}
      \vspace{0.2em}
      \quiet{Three details a referee asks about: the test is on
      \emph{probabilities}, not logits; the mask is \emph{detached}; it is
      recomputed \emph{every step}, never cached.}
    \end{column}
  \end{columns}
  \note{%
    AS-Softmax, published in 2023, is the method I set out to extend. Its rule
    is one line: for a training example, if the correct class already beats a
    competitor by at least delta in \emph{probability}, delete that competitor
    from the loss for this step. What remains is ordinary cross-entropy
    renormalised over the survivors. The picture: the true class is at 0.50,
    delta is 0.15, so the threshold sits at 0.35. The first competitor at 0.40
    is above the line --- still a genuine threat --- so it stays and keeps
    producing gradient. The two small ones are far below and get dropped. Three
    implementation points that a reviewer would ask about: the test is on
    probabilities, not on raw logits; the mask is detached, so no gradient flows
    through the threshold itself; and it is recomputed every single step, never
    cached. The published paper reports better accuracy, better calibration and
    about a 1.2 times speed-up.%
  }
\end{frame}

\begin{frame}{The gap in AS-Softmax --- and my hypothesis}
  \small
  \begin{alertblock}{$\dlt$ is a single scalar}
    The same number for ``sport vs.\ politics'' and ``sport vs.\ hardware''.
    The same number for a clean example and a mislabelled one.
    The same number at step $1$ and at step $10{,}000$.
  \end{alertblock}
  \vspace{-0.2em}
  \[
    \dlt \quad\longrightarrow\quad \mhit{\dltf}
    \qquad
    \tau = \frac{\text{current step}}{\text{total steps}} \in [0,1]
  \]
  \begin{columns}[T]
    \begin{column}{0.53\textwidth}
      \centering\scriptsize
      \begin{tabular}{@{}p{1.6cm}p{4.4cm}@{}}
        \toprule
        axis & why it might need its own margin \\
        \midrule
        class-pair $(t,j)$ & confusable pairs need a bigger gap than obvious ones \\
        sample $x$ & an ambiguous or mislabelled example is not like a clean one \\
        time $\tau$ & early training is noisy; start loose, tighten later \\
        \bottomrule
      \end{tabular}
    \end{column}
    \begin{column}{0.44\textwidth}
      \centering\scriptsize
      \begin{tabular}{@{}l>{\raggedright\arraybackslash}p{4.0cm}@{}}
        \toprule
        & hypothesis \\
        \midrule
        \textbf{H0} & AS-Softmax beats plain cross-entropy \quiet{(assumed, never tested)} \\
        \textbf{H1} & class-pair structure helps \\
        \textbf{H2} & per-sample structure helps \\
        \textbf{H3} & the axes compose \\
        \bottomrule
      \end{tabular}
    \end{column}
  \end{columns}
  \vspace{0.3em}
  \takeaway{\textbf{Success criterion, fixed before running anything:}
  $\geq +0.3\,\%$ accuracy over AS-Softmax, averaged over seeds. Writing it down
  in advance is what stops you moving the goalposts afterwards.}
  \note{%
    Here is the gap I went after. Delta is one scalar. It is the same number for
    a genuinely confusable pair of classes and for two classes that are nothing
    alike; the same for a clean example and a mislabelled one; the same at the
    first step and the last. That is a strong assumption and nobody had tested
    it. My proposal is to replace the scalar by a function of three arguments:
    the class pair, the sample, and tau --- normalised training progress between
    zero and one. That gives three hypotheses, H1 to H3. Note H0 at the top: the
    assumption that the method I am extending actually beats plain
    cross-entropy in the first place. I wrote it down as an assumption, not as
    something to test. That turns out to be the crux of the whole project.
    Finally, note the success criterion --- plus 0.3 percent over AS-Softmax,
    averaged over seeds --- fixed in advance, before any experiment ran, so I
    could not move the goalposts afterwards.%
  }
\end{frame}

% ===========================================================================
\section{Methodology}
% ===========================================================================

\begin{frame}{GAM-Softmax: the general form}
  \small
  \[
    \keep(x) \;=\; \{t\} \cup \bigl\{\, j \neq t \;:\; p_t - p_j < \dltf \,\bigr\},
    \qquad
    L \;=\; -\log\frac{e^{z_t}}{\displaystyle\sum_{k\in\keep(x)} e^{z_k}}
  \]
  Same loss as before --- the \emph{only} change is that the threshold is now a
  function rather than a constant.
  \begin{columns}[T]
    \begin{column}{0.52\textwidth}
      \begin{block}{Reduction property (unit-tested)}
        If $\dltf \equiv \dlt$ is constant, then $\keep$ is \emph{exactly}
        AS-Softmax's kept set, so $L \equiv L_{\mathrm{AS}}$ --- not
        approximately, identically.
      \end{block}
      \vspace{-0.2em}
      \begin{block}{Two familiar losses fall out}
        $\dlt \geq 1 \Rightarrow$ nothing is ever masked (since
        $p_t - p_j \leq 1$) $\Rightarrow$ \textbf{plain cross-entropy}.\\
        $\dlt$ constant $\Rightarrow$ \textbf{AS-Softmax}.
      \end{block}
    \end{column}
    \begin{column}{0.45\textwidth}
      \takeaway{This matters for honesty, not elegance: because the general form
      \emph{contains} the baseline, any difference in the results is
      attributable to the margin function alone --- not to a different loss,
      a different optimiser, or a different bug.}
      \vspace{0.3em}
      \quiet{\scriptsize Enforced by a test that asserts equality of the two
      losses to floating-point tolerance. 70 unit tests, $\approx 3$\,s.}
    \end{column}
  \end{columns}
  \note{%
    The generalisation is deliberately minimal. Same masking rule, same
    renormalised loss --- the single change is that the threshold on the right
    of the inequality is now a function of the class pair, the sample and the
    time, instead of a constant. Two consequences worth stating. If the function
    is constant, the kept set is exactly AS-Softmax's kept set and the losses are
    identical --- I assert that with a unit test, so it is checked on every
    commit. And if delta is at least one, nothing can ever be masked, because a
    difference of two probabilities can never exceed one --- so you recover plain
    cross-entropy. Both baselines are special cases of my formulation. That
    matters for scientific honesty rather than elegance: because the general
    form contains the baseline exactly, any difference in the numbers has to
    come from the margin function itself.%
  }
\end{frame}

\begin{frame}{M3 --- the class-pair axis}
  \small
  \[
    \dlt_{t,j} \;=\; \dlt_{\min} \;+\; \bigl(\dlt_{\max}(\tau) - \dlt_{\min}\bigr)\cdot
    \sigma\!\left(\frac{u_t \cdot v_j}{\sqrt{k}}\right),
    \qquad u, v \in \R^{n\times k},\ \ k = 8
  \]
  \begin{columns}[T]
    \begin{column}{0.52\textwidth}
      \textbf{Reading the formula}
      \begin{itemize}\setlength\itemsep{0.2em}
        \item every class $c$ gets two learnable vectors $u_c, v_c$ of length
              $k=8$; their inner product scores the \emph{ordered} pair $(t,j)$
        \item $\sigma$ (the logistic sigmoid) squashes that score into $(0,1)$,
              so $\dlt_{t,j}$ can never leave $[\dlt_{\min}, \dlt_{\max}]$
        \item $\dlt_{\max}(\tau)$ is the \emph{time} axis: a warm-up schedule
        \item dividing by $\sqrt{k}$ keeps the pre-sigmoid scale $O(1)$
              --- the same variance argument as in attention
      \end{itemize}
    \end{column}
    \begin{column}{0.45\textwidth}
      \begin{block}{Why low-rank, not a full table?}
        A free margin per ordered pair is $n(n-1) = 380$ unconstrained numbers
        for $n=20$, with no sharing between classes.\\[0.3em]
        The low-rank form uses $2nk = 320$ and \emph{forces} classes to share
        structure --- which is the actual hypothesis: that pair difficulty is
        low-dimensional.
      \end{block}
    \end{column}
  \end{columns}
  \note{%
    The first axis. Each class gets two short learnable vectors, u and v, of
    length eight. The margin for the ordered pair --- true class t against
    competitor j --- is the inner product of u-t and v-j, pushed through a
    sigmoid so it is bounded, and then stretched between a floor and a
    time-varying ceiling. The square-root-of-k divisor is the same variance
    normalisation used in attention: it stops the inner product from saturating
    the sigmoid at initialisation. Why low-rank rather than a free table? For
    twenty classes a free table is 380 unconstrained numbers with no sharing at
    all, and it would just overfit. The factored form is 320 parameters that
    force classes to share structure --- and that sharing \emph{is} the
    hypothesis: that how confusable a pair is, is a low-dimensional property.%
  }
\end{frame}

\begin{frame}{M4 --- the sample axis}
  \small
  \textbf{Step 1 --- how suspicious is this example, relative to its batch?}
  \[
    s_i \;=\; \sigma\!\left(\frac{\bar{p}_t - p_{t,i}}{\operatorname{std}(p_t)\cdot\text{temp}}\right)
    \in (0,1),
    \qquad s_i > 0.5 \iff \text{less confident than its peers}
  \]
  \textbf{Step 2 --- move the margin by that amount}
  \[
    \dlt_i(\tau) \;=\; \operatorname{clamp}\Bigl(\underbrace{\dlt_{\text{base}}(\tau)}_{\text{time axis}}
    \;-\; \beta(\tau)\,\bigl(2 s_i - 1\bigr),\ \ \dlt_{\text{floor}},\ \dlt_{\text{ceil}}\Bigr)
  \]
  \begin{columns}[T]
    \begin{column}{0.52\textwidth}
      \begin{itemize}\setlength\itemsep{0.2em}
        \item suspicious ($s_i\!\to\!1$) $\Rightarrow$ $\dlt$ \emph{down}
              $\Rightarrow$ more classes masked $\Rightarrow$ less gradient
        \item confident ($s_i\!\to\!0$) $\Rightarrow$ $\dlt$ \emph{up}, but
              capped at AS-Softmax's own $\dlt$ --- the axis is \alert{one-sided},
              so any difference comes purely from what we do to suspicious samples
      \end{itemize}
    \end{column}
    \begin{column}{0.45\textwidth}
      \begin{block}{Two deliberate design choices}
        $(2s_i-1)$ averages to $\approx 0$ over a batch, so this varies $\dlt$
        \emph{across} samples without changing the overall masking level ---
        that stays the time axis's job, keeping the axes separable.\\[0.3em]
        $\beta(\tau)=0$ for the first 30\,\% of training: early on $p_t$ is
        near-random and carries no signal about which labels are wrong.
      \end{block}
    \end{column}
  \end{columns}
  \note{%
    The second axis, and this is the one that worked. Two steps. First, score
    how suspicious each example is, using only information the model already
    computed: take the probability it currently assigns to that example's
    training label, compare it to the batch mean, divide by the batch standard
    deviation, and squash with a sigmoid. Above one half means ``this example is
    less confident than its peers'' --- that is, suspicious. Second, shift the
    margin down by that amount. Suspicious examples get a lower margin, so more
    of their competitors get masked, so they contribute less gradient. Two
    design choices I would defend to a referee. The correction is mean-centred,
    so it redistributes the margin across samples without changing the average
    masking level --- that keeps this axis independent of the time axis, which
    matters for the ablation. And beta is zero for the first thirty percent of
    training, because early on the model's confidence is near-random and tells
    you nothing about which labels are wrong.%
  }
\end{frame}

\begin{frame}{The idea worth keeping: let $\dlt$ go \emph{negative}}
  \small
  Nothing in the mathematics requires a margin to be positive. Since
  $p_t - p_j \in [-1, 1]$, the value $\dlt = -0.85$ is a perfectly legal
  threshold. It reads:
  \begin{center}
    \emph{``drop class $j$ even if it is currently beating the training label by up to $0.85$.''}
  \end{center}
  \begin{columns}[T]
    \begin{column}{0.49\textwidth}
      \textbf{On a mislabelled example}
      \begin{itemize}\setlength\itemsep{0.15em}
        \item the model keeps $p_t$ low (the label is wrong)
        \item so with $\dlt > 0$, \emph{no} competitor ever clears the bar
        \item nothing is masked $\Rightarrow$ full cross-entropy $\Rightarrow$
              pushed to fit the wrong label forever
        \item \alert{that is exactly how a network memorises noise}
      \end{itemize}
    \end{column}
    \begin{column}{0.49\textwidth}
      \textbf{With $\dlt < 0$, the algebra ends it}
      \[ \keep(x) = \{t\} \ \Rightarrow\ p_t^{\text{renorm}} = \frac{e^{z_t}}{e^{z_t}} = 1 \]
      \[ \Rightarrow\ L = -\log 1 = 0 \ \Rightarrow\ \nabla L = 0 \]
      The suspicious sample \emph{removes itself} from training --- no threshold
      rule, no sample-selection heuristic.
    \end{column}
  \end{columns}
  \begin{alertblock}{The reframing}
    The noisy-label literature's ``small-loss trick'' --- drop the
    high-loss examples --- is exactly the special case $\dlt < 0$ of a
    \emph{margin}. No second network, no auxiliary loss, no selection
    schedule: one line of configuration.
  \end{alertblock}
  \note{%
    This is the slide I would keep if I could keep only one. Nothing in the maths
    says a margin has to be positive. The quantity being thresholded is a
    difference of two probabilities, so it lives between minus one and plus one,
    and minus 0.85 is a legal threshold. It means: drop this competitor even if
    it is currently \emph{ahead} of the training label. Now follow a mislabelled
    example. Because the label is wrong, the model keeps assigning it low
    probability. So under a positive margin no competitor ever clears the bar,
    nothing gets masked, and the example receives full cross-entropy pressure
    forever --- that is the memorisation mechanism, precisely. Under a negative
    margin, every competitor is masked, the kept set is the single true class,
    the renormalised probability is e-to-the-z over e-to-the-z which is exactly
    one, the log of one is zero, and the gradient is zero. The example removes
    itself. And that gives the reframing in the box: dropping high-loss samples
    --- a standard trick in the noisy-label literature that usually needs a
    second network or a selection schedule --- is just a margin with a negative
    delta. One config line.%
  }
\end{frame}

% ===========================================================================
\section{Experimentation}
% ===========================================================================

\begin{frame}{The setup}
  \small
  \begin{columns}[T]
    \begin{column}{0.5\textwidth}
      \textbf{Constraint that shaped everything}\\
      One \textbf{RTX 3050 laptop, 6\,GB} (the plan assumed a 3090/4090).
      One 20-Newsgroups run $\approx 25$ min $\Rightarrow$ a properly seeded
      sweep is a day of wall clock.
      \vspace{0.5em}

      \textbf{Testbed}
      \begin{itemize}\setlength\itemsep{0.1em}
        \item 20 Newsgroups, 20 classes, headers/footers/quotes \emph{stripped}
              --- otherwise BERT reaches $\approx99\,\%$ from metadata and no
              method can separate from another
        \item BERT-base, 4 epochs, batch 16, lr $2\!\times\!10^{-5}$, seq len 128
        \item 6 configs, \textbf{identical} except the \texttt{loss:} block; 2 seeds (42, 43)
      \end{itemize}
    \end{column}
    \begin{column}{0.47\textwidth}
      \begin{block}{Label noise --- and why it is a \emph{paired} comparison}
        \textbf{40\,\% symmetric label noise} on the training set; validation
        labels stay clean.\\[0.3em]
        The noise seed is fixed and \emph{decoupled} from the run seed, so every
        method sees the \emph{identical} corrupted labels. Methods are compared
        on the same corrupted dataset, not on statistically similar ones.
      \end{block}
      \takeaway{\textbf{One config file = one experiment = one row of the results
      table.} 9 loss functions, 70 unit tests running in $\approx3$\,s. The null
      result is not a code failure.}
    \end{column}
  \end{columns}
  \note{%
    The hardware constraint is real and it shaped every choice: a six-gigabyte
    laptop GPU, where one run of this dataset takes about twenty-five minutes,
    so a properly seeded comparison is a full day. The testbed is 20 Newsgroups
    with headers, footers and quotes stripped --- without stripping, BERT scores
    about ninety-nine percent by reading email metadata and every method ties, so
    the experiment would measure nothing. Then the key protocol point on the
    right: I corrupt forty percent of the training labels, at random, to a
    different class, while validation labels stay clean. The noise seed is fixed
    and separate from the run seed, which means every method sees the exact same
    corrupted labels. So this is a paired comparison --- the methods differ, the
    data does not. Finally the engineering discipline: one config file is one
    experiment is one row of the results table, and seventy unit tests pass in
    about three seconds. When I show you a null result, it is not a bug.%
  }
\end{frame}

\begin{frame}{What I measure --- and the arithmetic behind it}
  \small
  Accuracy alone had been flat in every earlier experiment, so E1 measures the
  \emph{mechanism} directly.
  \begin{columns}[T]
    \begin{column}{0.52\textwidth}
      \begin{itemize}\setlength\itemsep{0.25em}
        \item \textbf{peak} and \textbf{final} validation accuracy, and the
              \textbf{decay} between them --- the overfitting trajectory
        \item \textbf{memorisation rate:} of the training examples whose labels
              were \emph{corrupted}, the fraction the model now predicts
              \emph{as their corrupted label}
        \item \textbf{recovery rate:} same examples, fraction predicted as their
              \emph{true} label despite being trained on a wrong one
        \item \textbf{masked ratio}, and 15-bin \textbf{ECE} (calibration)
      \end{itemize}
    \end{column}
    \begin{column}{0.45\textwidth}
      \begin{block}{The no-memorisation floor}
        With 40\,\% noise, where a flip never returns the original label, a model
        that learned the \emph{true} function and memorised \emph{nothing}
        scores exactly
        \[ 0.6 \times 100\,\% + 0.4 \times 0\,\% = \mhit{60.0\,\%} \]
        on the labels it was given.\\[0.3em]
        \textbf{Everything above 60 is memorised noise} --- a quantity, not an
        impression.
      \end{block}
    \end{column}
  \end{columns}
  \note{%
    Six weeks of reading only accuracy taught me nothing, because accuracy was
    flat everywhere. So the final experiment measures the mechanism. The central
    quantity is the memorisation rate: I keep a record of exactly which training
    examples had their labels corrupted, and at the end I ask what fraction of
    those the model now predicts as the wrong label it was trained on. That is
    memorisation measured directly, not inferred from an accuracy curve. And the
    box on the right is the arithmetic that makes it interpretable. Sixty percent
    of the training labels are still correct; forty percent were flipped, and a
    flip never returns the original label. So a model that perfectly learned the
    true function and memorised nothing would score exactly sixty percent
    against the labels it was handed. Everything above sixty is memorised noise.
    That gives a principled zero point.%
  }
\end{frame}

\begin{frame}{Every experiment, in order}
  \scriptsize
  \begin{table}\centering
  \begin{tabular}{@{}clp{2.5cm}p{5.0cm}@{}}
    \toprule
    \# & testbed & question & result \\
    \midrule
    1 & SST-5, clean, 3 seeds & baselines sane? & softmax 51.1\,\%, AS-Softmax \textbf{52.2\,\%} $\Rightarrow$ AS $>$ CE by $+1.1$ \checkmark \\
    2 & SST-5, clean, 3 seeds & H1, fixed margin & M3 52.10 vs AS 52.17 $=$ \hit{$-0.06$} \\
    3 & SST-5, clean, 3 seeds & H1, learnable margin & M3$'$ 51.83 vs AS 52.17 $=$ \hit{$-0.33$} \\
    4 & SST-5, tuning screen & is H1 under-tuned? & higher margin LR hurt \emph{monotonically} ($52.5 \to 50.6 \to 50.1$) $\Rightarrow$ not a tuning artefact \\
    5 & 20NG, clean, 1 seed & testbed with headroom? & GAM 71.18 $>$ AS 70.59, but plain CE 71.46 wins $\Rightarrow$ \hit{AS $<$ CE} \\
    6 & SST-5, noise $\{0,.2,.4\}$ & does a hard regime rescue it? & all gaps inside $\pm2\,\%$ seed scatter $\Rightarrow$ inconclusive \\
    \midrule
    \textbf{7} & \textbf{20NG, 40\,\% noise, 2 seeds} & \textbf{the decisive test} & \textbf{next section} \\
    \bottomrule
  \end{tabular}
  \end{table}
  \begin{columns}[T]
    \begin{column}{0.52\textwidth}
      \takeaway{\textbf{H0 kept failing.} AS-Softmax beat cross-entropy on
      SST-5 ($+1.1$), lost on 20NG ($-0.9$), and was a wash under noise. Three
      of four setups gave the method I extend \emph{no edge} over the simplest
      possible baseline.}
    \end{column}
    \begin{column}{0.45\textwidth}
      \takeaway{\textbf{The one repeatable GAM signal was stability, not
      accuracy}: lowest seed-to-seed variance at every noise level, and still
      \emph{improving} at the last epoch where the baselines had peaked and
      turned over.}
    \end{column}
  \end{columns}
  \note{%
    The honest audit trail --- every experiment, in the order I ran them.
    Row one: baselines reproduce, AS-Softmax beats plain softmax by about one
    point on SST-5, so the engine is trustworthy. Rows two and three: the
    class-pair hypothesis, tested twice, comes out slightly negative both times.
    Row four is the one I would want a referee to see --- I checked whether the
    flat result was just bad tuning, and raising the margin learning rate made
    it monotonically \emph{worse}, so it is a real effect, not under-tuning.
    Row five moves to a harder dataset and finds the thing that eventually
    decided the project: plain cross-entropy wins, and AS-Softmax --- the method
    I am extending --- loses to it. Row six, adding label noise, is
    inconclusive. Two patterns repeat. First, H0 --- my unexamined assumption ---
    kept failing. Second, the only repeatable signal from my method was
    stability rather than accuracy, and that is what the final experiment was
    designed to confirm or kill.%
  }
\end{frame}

% ===========================================================================
\section{Results and Demonstration of Final Result}
% ===========================================================================

\begin{frame}{E1 --- the decisive experiment}
  \framesubtitle{20 Newsgroups, 40\,\% symmetric label noise, BERT-base, 2 seeds (42, 43)}
  \scriptsize
  \begin{table}\centering
  \begin{tabular}{@{}lrrrrrr@{}}
    \toprule
    method & peak val \% & final val \% & decay & \multicolumn{1}{c}{\shortstack{memorised noise\\(pts over floor)}} & mem @final \% & recovers true \% \\
    \midrule
    softmax (plain CE)        & 67.59 {\tiny$\pm$0.71} & 64.19 {\tiny$\pm$0.92} & $-3.40$ & $+11.5$ & 33.7 {\tiny$\pm$0.4} & 41.2 \\
    as\_softmax (scalar $\dlt$) & 67.56 {\tiny$\pm$0.34} & 64.72 {\tiny$\pm$0.20} & $-2.84$ & \hit{$+13.6$} & \hit{39.9 {\tiny$\pm$1.8}} & 44.1 \\
    gam\_m3 (class-pair)      & 67.68 {\tiny$\pm$0.81} & 65.21 {\tiny$\pm$1.42} & $-2.47$ & $+11.4$ & 35.2 {\tiny$\pm$1.7} & 45.9 \\
    \textbf{gam\_m4 (sample)} & \textbf{68.09} {\tiny$\pm$0.05} & \good{66.96} {\tiny$\pm$0.42} & \good{$-1.13$} & \good{$+3.1$} & \good{19.3} {\tiny$\pm$0.7} & \good{57.3} \\
    \bottomrule
  \end{tabular}
  \end{table}
  \vspace{-0.3em}
  \begin{columns}[T]
    \begin{column}{0.49\textwidth}
      \begin{block}{Read the accuracy column first --- it is flat}
        Peak accuracy spans $67.56$ to $68.09$: half a point across all four
        methods. Plain cross-entropy alone spans \textbf{1.0 point} between its
        two seeds. \textbf{The accuracy column contains no result.}
      \end{block}
    \end{column}
    \begin{column}{0.49\textwidth}
      \begin{alertblock}{Now read the memorisation column}
        The same two seeds of plain cross-entropy span only \textbf{0.6 points}
        of memorisation. So a \textbf{14-point} gap is $\approx20\times$ the
        measurement's own noise --- it cannot be a lucky seed.
      \end{alertblock}
    \end{column}
  \end{columns}
  \note{%
    Here is the whole experiment on one slide. Read the accuracy column first,
    because the discipline matters more than the result: peak accuracy runs from
    67.6 to 68.1, about half a point across four quite different objectives.
    Meanwhile plain cross-entropy on its own spans a full point between its two
    random seeds. So the accuracy column contains no result --- every difference
    in it is smaller than the noise of the measurement. Now the memorisation
    column. The same two seeds of cross-entropy span only six tenths of a point
    there, so that measurement is roughly twenty times quieter. And the gap
    between the bottom row and the rest is fourteen points. That is about twenty
    times the noise floor of the measurement, which is why I am willing to make
    a claim about that column and not about the accuracy column. Choosing which
    metric can support a claim is most of this experiment.%
  }
\end{frame}

\begin{frame}{The output, in two pictures}
  \begin{columns}[T]
    \begin{column}{0.5\textwidth}
      \centering
      \includegraphics[width=0.93\linewidth]{fig2_memorization_n0.4.png}
      \vspace{-0.3em}
      {\scriptsize Lower is better. All methods rise --- \alert{M4's curve
      flattens} while the rest accelerate.}
    \end{column}
    \begin{column}{0.5\textwidth}
      \centering
      \includegraphics[width=0.93\linewidth]{fig1_accuracy_n0.4.png}
      \vspace{-0.3em}
      {\scriptsize Everything peaks at epoch 1 and decays. \alert{M4 decays
      least} --- $-1.13$ points against cross-entropy's $-3.40$.}
    \end{column}
  \end{columns}
  \vspace{0.2em}
  \takeaway{Same runs, two views. Left: by the last epoch, AS-Softmax has
  memorised \emph{more} corrupted labels than plain cross-entropy (39.9\,\% vs
  33.7\,\%) --- the opposite of what the margin-loss story predicts. Right: the
  four curves are almost on top of each other until the overfitting phase
  separates them.}
  \note{%
    The two pictures are the same four runs, viewed two ways. On the left,
    memorisation of corrupted labels over the four epochs. Everything climbs ---
    all of these models memorise --- but the red curve, the sample-axis variant,
    flattens out around fifteen to nineteen percent while the others accelerate
    past thirty-five. Note the blue curve, AS-Softmax: it ends \emph{highest},
    above plain cross-entropy. I will come back to that, because it contradicts
    the standard justification for these losses. On the right, validation
    accuracy on clean labels. Everything peaks at epoch one and then decays as
    the models start fitting the noise. The red curve decays least --- about one
    point lost against cross-entropy's three and a half. Same mechanism, two
    views: less memorisation, less decay.%
  }
\end{frame}

\begin{frame}{What the numbers actually say}
  \small
  \begin{columns}[T]
    \begin{column}{0.49\textwidth}
      \textbf{1. The precondition failed.}\\
      AS $-$ CE $= -0.03$ points over two seeds ($-0.29$ and $+0.23$: a coin
      flip). \emph{Fourth testbed in a row.} The method I set out to generalise
      has, in my hands, no advantage over plain cross-entropy to generalise.
      \vspace{0.5em}

      \textbf{2. AS-Softmax memorises \emph{more} than plain CE.}\\
      $+13.6$ vs $+11.5$ points over the floor, on \emph{both} seeds
      ($+7.8$, $+4.6$). This contradicts the standard story --- and my own
      reason for moving to noisy labels.\\[0.3em]
      \quiet{Mechanism: masking retires the \emph{easy} negatives first --- what
      a correctly-labelled example runs out of early. So an ever-larger share of
      the gradient goes to the still-confusable examples, which are exactly the
      mislabelled ones. \textbf{Scalar masking concentrates effort onto the
      noise.}}
    \end{column}
    \begin{column}{0.49\textwidth}
      \textbf{3. The sample axis does what it was designed to do.}
      \vspace{0.2em}
      \begin{table}\centering\scriptsize
      \begin{tabular}{@{}lrrrr@{}}
        \toprule
        (2-seed means) & CE & AS & M3 & \textbf{M4} \\
        \midrule
        memorised noise (pts) & 11.5 & 13.6 & 11.4 & \good{3.1} \\
        recovers true label & 41.2 & 44.1 & 45.9 & \good{57.3} \\
        peak-to-final decay & 3.40 & 2.84 & 2.47 & \good{1.13} \\
        non-target slots masked & 0\,\% & 56\,\% & 49\,\% & \good{95\,\%} \\
        \bottomrule
      \end{tabular}
      \end{table}
      Four independent measurements moving together in the predicted direction,
      driven by a mechanism I can point at in the code: $\dlt$ goes negative
      $\to$ the suspect sample's competitors are all masked $\to$ its loss
      contribution vanishes.
    \end{column}
  \end{columns}
  \note{%
    Three readings. First, the precondition. Over two seeds AS-Softmax beats
    plain cross-entropy by minus 0.03 points --- minus 0.29 on one seed, plus
    0.23 on the other, which is a coin flip. That is the fourth testbed in a row
    where the method I set out to generalise has no reliable edge over the
    simplest possible baseline. You cannot generalise an advantage that is not
    there. Second, and this one surprised me: AS-Softmax memorises corrupted
    labels \emph{more} than plain cross-entropy, on both seeds. Every
    margin-loss paper, and my own written rationale for moving to noisy labels,
    argues that masking should \emph{prevent} memorisation. Measured directly,
    the opposite happens, and I think I know why --- masking retires the easy
    competitors first, and easy competitors are what a correctly labelled
    example runs out of early, so as training goes on an ever-larger share of
    the gradient budget is spent on the hard examples, which are precisely the
    mislabelled ones. Scalar masking concentrates effort onto the noise. Third,
    the table: four independent measurements all move together for the
    sample-axis variant, and I can point at the line of code that causes it.%
  }
\end{frame}

\begin{frame}{Verdict --- stated as it is}
  \scriptsize
  \begin{columns}[T]
    \begin{column}{0.52\textwidth}
      \begin{table}\centering
      \begin{tabular}{@{}llp{2.6cm}@{}}
        \toprule
        & hypothesis & verdict \\
        \midrule
        \textbf{H0} & AS $>$ plain CE & \hit{NOT SUPPORTED} --- 3 of 4 testbeds show no advantage \\
        \textbf{H1} & class-pair helps & \hit{NOT SUPPORTED} --- $+0.12$ over AS, inside a 1.0-pt noise band \\
        \textbf{H2} & sample axis helps & \textbf{SPLIT} --- peak accuracy not established; memorisation cut $\approx4\times$ \\
        \textbf{H3} & the axes compose & UNTESTED --- implemented, run cut for compute \\
        \bottomrule
      \end{tabular}
      \end{table}
      \vspace{0.2em}
      \begin{alertblock}{The honest limit}
        M4's mean peak is $+0.51$ over CE --- which would clear my $+0.3$
        criterion. \textbf{I do not claim it.} Per seed it is $-0.03$ and
        $+1.04$: the entire mean comes from one seed where cross-entropy had a
        bad run. What \emph{does} replicate is the \emph{final-epoch} gain:
        $+3.72$ and $+1.83$, both positive, mean $+2.77$.
      \end{alertblock}
    \end{column}
    \begin{column}{0.46\textwidth}
      \begin{table}\centering
      \begin{tabular}{@{}p{2.4cm}p{3.4cm}@{}}
        \toprule
        finding & confidence \\
        \midrule
        M4 cuts memorisation $\approx4\times$ ($-14.4$ pts) & \good{Confirmed} \quiet{(both seeds; $\approx20\times$ the metric's spread)} \\
        AS-Softmax memorises \emph{more} than CE & \good{Confirmed} \quiet{(both seeds)} \\
        M4 gains at the final epoch ($+2.77$) & \good{Confirmed (direction)} \quiet{(magnitude varies $2\times$)} \\
        H0 and H1 null & \good{Confirmed} \quiet{(4 testbeds)} \\
        M4 gains at \emph{peak} accuracy & \hit{Not established} \quiet{(one seed)} \\
        M4 is more seed-stable & \quiet{Suggestive only ($n=2$)} \\
        Calibration & \quiet{No claim} \\
        \bottomrule
      \end{tabular}
      \end{table}
    \end{column}
  \end{columns}
  \note{%
    The verdict table, left. My assumed precondition: not supported. The
    class-pair hypothesis, which was the original core of the project: not
    supported, and I showed earlier it is not a tuning artefact. The sample axis:
    split --- it fails on the metric I pre-registered and succeeds on the
    mechanism metrics. The third hypothesis was implemented and unit-tested but I
    could not afford the compute to run it, and I would rather say that than
    quietly omit it. Then the box, which is the part I want to be judged on.
    The sample-axis variant's mean peak accuracy is plus 0.51 over cross-entropy,
    which would clear the plus 0.3 criterion I set in advance. I do not claim it.
    Per seed the gap is minus 0.03 and plus 1.04 --- the entire mean comes from
    one seed where cross-entropy happened to have a bad run, and one seed of two
    is not a result. What does replicate is the end-of-training gain, positive on
    both seeds. On the right, every finding graded separately, because they are
    not equally solid.%
  }
\end{frame}

\begin{frame}{Conclusion}
  \small
  \begin{block}{In one sentence}
    Generalising AS-Softmax's scalar margin along the \textbf{class-pair} axis
    produced no measurable benefit on any of four testbeds; generalising it
    along the \textbf{sample} axis --- by letting the margin go
    \emph{negative}, which makes small-loss sample rejection a special case of a
    margin --- did not reliably improve peak accuracy either, but \textbf{cut
    memorisation of corrupted labels roughly fourfold} and left the final-epoch
    model \textbf{2.8 points} better than cross-entropy under 40\,\% label noise.
  \end{block}
  \begin{columns}[T]
    \begin{column}{0.485\textwidth}
      \scriptsize
      \textbf{What I would claim in writing}
      \begin{itemize}\setlength\itemsep{0.15em}
        \item the class-pair negative result, replicated over four testbeds and
              shown not to be a tuning artefact
        \item that AS-Softmax \emph{increases} memorisation under label noise,
              contradicting the standard rationale for margin losses in that
              regime --- measured directly, not inferred from accuracy
        \item the negative-margin reframing (small-loss rejection \emph{is} a
              margin), with evidence that it works mechanically
      \end{itemize}
    \end{column}
    \begin{column}{0.485\textwidth}
      \scriptsize
      \textbf{What I would \emph{not} claim}
      \begin{itemize}\setlength\itemsep{0.15em}
        \item that GAM-Softmax improves accuracy. It does not.
        \item any state-of-the-art noisy-label claim --- I never ran the real
              baselines (co-teaching, DivideMix, GCE). M4 beats \emph{cross-entropy
              and AS-Softmax at the final epoch}; that is a much weaker statement,
              and it is the only one the data supports.
        \item any calibration claim \quiet{(the ECE gap is \emph{under}confidence;
              largely fixable by temperature scaling, which I did not run)}
      \end{itemize}
    \end{column}
  \end{columns}
  \note{%
    The one-sentence conclusion is in the box, and it is written so that nothing
    in it exceeds the evidence. Then the two lists. On the left, what I would
    defend in writing: the negative result on the class-pair axis, replicated
    across four testbeds and shown not to be a tuning artefact; the finding that
    AS-Softmax increases memorisation, which contradicts the standard rationale
    and which I measured directly rather than inferred; and the negative-margin
    reframing as a modelling contribution. On the right, what I refuse to claim
    --- above all that the method improves accuracy, because it does not.%
  }
\end{frame}

\begin{frame}{Limits, and what a follow-up should do differently}
  \scriptsize
  \begin{columns}[T]
    \begin{column}{0.485\textwidth}
      \textbf{What would have to be true for this to be wrong}
      \begin{itemize}\setlength\itemsep{0.2em}
        \item 2 seeds, 1 dataset, 1 noise level, 1 architecture, 4 epochs.
              \quiet{Two seeds is enough to kill a hypothesis and not enough to
              establish a small effect --- which is exactly why the memorisation
              gap is claimed and the accuracy gap is not.}
        \item $\beta$ was never swept --- the cleanest evidence for a mechanism
              is that the effect scales with its knob, and that is the run I
              could not afford
        \item symmetric noise is the easy case; real noise is
              instance-dependent, where ``the model disagrees with this label''
              is a much weaker signal of corruption
        \item weak baselines by design --- CE and AS-Softmax, not methods built
              for noisy labels
      \end{itemize}
    \end{column}
    \begin{column}{0.485\textwidth}
      \textbf{What a follow-up should do differently}
      \begin{enumerate}\setlength\itemsep{0.2em}
        \item \alert{validate H0 in week 1}, not week 8. One week reproducing
              the baseline's published gain would have redirected the whole
              project.
        \item pick a regime where masking is \emph{forced} to matter --- 10k+
              classes, or genuinely long-tailed data. Balanced text
              classification with a strong pretrained encoder is close to the
              worst case: plain CE is very hard to beat there.
        \item instrument the mechanism from day 1. I read only accuracy for six
              weeks; it was flat everywhere and told me nothing about \emph{why}.
        \item keep the negative-margin framing, and test it against the real
              noisy-label baselines. It is one config line away.
      \end{enumerate}
    \end{column}
  \end{columns}
  \note{%
    The limits, stated before anyone has to ask. Two seeds, one dataset, one
    noise level --- enough to kill a hypothesis, not enough to establish a small
    effect, which is precisely the distinction I drew two slides ago. Beta was
    never swept, and that is the honest hole: the cleanest evidence that a
    mechanism is what you think it is, is that the effect scales with its knob,
    and that is the run I could not afford. Symmetric noise is the easy case,
    and the baselines are weak by design. On the right, the four things I would
    do differently, and the first is the real lesson of the whole project: I
    assumed the baseline's advantage instead of spending one week checking it,
    and that assumption was the crux. Thank you --- happy to take questions.%
  }
\end{frame}

% =========================================================== BACKUP
\appendix

\begin{frame}[plain]
  \vfill\centering
  {\color{gamgrey}\small BACKUP SLIDES}\\[0.6em]
  {\color{gamblue}\Large\bfseries For questions}
  \vfill
  \note{Backup slides follow: per-seed numbers, the calibration caveat, and how to reproduce everything.}
\end{frame}

\begin{frame}{Backup --- head-to-head, reported per seed}
  \small
  A two-seed mean can hide the fact that one seed did all the work, so here is
  every comparison split out.
  \begin{table}\centering\small
  \begin{tabular}{@{}lccc@{}}
    \toprule
    vs.\ plain CE & peak (s42 / s43) & final (s42 / s43) & memorisation (s42 / s43) \\
    \midrule
    as\_softmax $-$ CE & $-0.29$ / $+0.23$ & $+1.33$ / $-0.26$ & \hit{$+7.8$ / $+4.6$} \\
    gam\_m3 $-$ CE     & $+0.16$ / $+0.01$ & $+0.67$ / $+1.38$ & $+0.6$ / $+2.4$ \\
    gam\_m4 $-$ CE     & $-0.03$ / $\mathbf{+1.04}$ & \good{$+3.72$ / $+1.83$} & \good{$-14.6$ / $-14.2$} \\
    \bottomrule
  \end{tabular}
  \end{table}
  \takeaway{Read the M4 row across: the peak column disagrees with itself, the
  final and memorisation columns agree on both seeds. That difference is the
  entire basis for which claims I make and which I refuse to.}
  \note{%
    This is the table that decides what I am willing to say. Look at the bottom
    row. The peak-accuracy column disagrees with itself --- minus 0.03 on one
    seed, plus 1.04 on the other. The final-accuracy and memorisation columns
    agree on both seeds, and the memorisation effect is fourteen points in both
    cases. So I claim the second and third columns and refuse the first.%
  }
\end{frame}

\begin{frame}{Backup --- the calibration caveat}
  \small
  The margin variants have far \emph{worse} ECE: $0.42$--$0.47$ against plain
  cross-entropy's $0.19$. Read naively that is a defeat. It is not what is
  happening.
  \begin{columns}[T]
    \begin{column}{0.52\textwidth}
      \begin{itemize}\setlength\itemsep{0.2em}
        \item M4's mean predicted confidence is $\approx 0.25$ while its
              accuracy is $\approx 0.67$
        \item so the model is systematically \alert{under}confident, not
              overconfident
        \item by construction: the margin objective \emph{stops pushing} $p_t$
              upward once the gap is met, so the probabilities stay compressed
      \end{itemize}
    \end{column}
    \begin{column}{0.45\textwidth}
      \takeaway{Underconfidence is a benign failure mode compared with
      cross-entropy's direction, and it is largely removable by temperature
      scaling --- which I did not run. So: \textbf{no calibration claim in
      either direction.}}
    \end{column}
  \end{columns}
  \note{%
    If someone asks about calibration: expected calibration error is much worse
    for the margin variants, 0.42 to 0.47 against 0.19. But look at the
    direction. Mean predicted confidence for the sample-axis variant is about
    0.25 while its accuracy is about 0.67 --- so it is badly \emph{under}
    confident, not overconfident, and that is exactly what the objective is
    designed to do, because it stops pushing the true-class probability upward
    once the margin is met. Underconfidence is the benign direction and it is
    largely removable by temperature scaling, which I did not run. So I make no
    calibration claim either way.%
  }
\end{frame}

\begin{frame}[fragile]{Backup --- reproducing all of it}
  \small
  \begin{semiverbatim}\scriptsize
conda activate gam

python -m pytest tests/ -v                    \textcolor{gamgrey}{-- 70 tests, ~3 s}
python experiments/final_experiment.py        \textcolor{gamgrey}{-- E1, resumable (~6 h)}
python experiments/final_experiment.py --only-block 1
  \end{semiverbatim}
  \begin{columns}[T]
    \begin{column}{0.48\textwidth}
      \begin{itemize}\setlength\itemsep{0.1em}
        \item \texttt{FINAL\_REPORT.md} --- the full write-up
        \item \texttt{HOW\_WE\_VARY\_DELTA.md} --- the method in plain language
        \item \texttt{runs/final\_20ng/} --- per-run JSON with full epoch
              histories, plus \texttt{summary.md}
      \end{itemize}
    \end{column}
    \begin{column}{0.48\textwidth}
      \begin{itemize}\setlength\itemsep{0.1em}
        \item 9 losses in \texttt{gam\_softmax/losses/}
        \item margins in \texttt{gam\_softmax/margins/}
        \item one config = one experiment = one table row
        \item run 2026-07-28/31, commit \texttt{b2ff750}
      \end{itemize}
    \end{column}
  \end{columns}
  \note{%
    Everything is reproducible from the repository: the test suite runs in three
    seconds, and the final experiment is a single resumable command. Every
    number in this talk traces back to a JSON record under runs slash
    final-20ng, and the runs are stamped with the commit they came from.%
  }
\end{frame}
